\chapter{Classical Mechanics}
\label{cm}

{
\setlength\epigraphwidth{0.6\textwidth}
\epigraph {
    Motion\\
There is no motion, said a wise old Greek.\\
The other sage just smiled and walked around.\\
A better answer hardly could be found,\\
And everybody praised this tongue in cheek.\\
But, gentlemen, this funny story might\\
Remind us of a man of great renown:\\
Day after day the Sun walks up and down,\\
And yet the stubborn Galilei was right.
} {A.S. Pushkin, Translated by D. Broytman }
}

\section{The Two Roles of Classical Mechanics}
\label{cm:tr}

One of the most important models in physics is mechanics.
Until the end of the 19th century, mechanics was considered as direct and the only
the basis of all physics. By words understand or
to explain a physical phenomenon meant constructing a mechanical model of it.

Situation has changed,  other physical models exist, e.g. quantum mechanics so 
"usual" mechanics has changed its name to "classical mechanics" and still is
the central part of theoretical physics. All main physical concepts emerged from
it.

Classical mechanics occupies a peculiar position among physical theories. It is,
first, a complete and self-contained subject: a body of concepts, equations, and
results describing the motion of bodies under forces, valid and useful on its own
terms, with no reference whatsoever to the Planck constant $\hbar$ or the speed of
light $c$. Celestial mechanics, rigid-body dynamics, the theory of small
oscillations, continuum mechanics, and the modern theory of dynamical systems and
chaos are all classical mechanics pursued for its own sake, and none of them is
waiting to be corrected by a more fundamental theory in order to be meaningful.

Second,  classical mechanics is
the \emph{base} on which relativistic and quantum mechanics are built. Special and
general relativity are commonly introduced as corrections to Newtonian mechanics
that become significant when velocities approach $c$ or fields become strong;
quantum mechanics is commonly introduced as a modification that becomes significant
at the scale of $\hbar$. In both cases, the classical theory is treated as a known,
well-understood limit --- the solid ground from which the generalization departs
and to which it must reduce.

It is tempting to think that once this second, subordinate role is granted, nothing
further needs to be said about classical mechanics: its job is simply to sit still
and be a limit. Here we argue that this is false in two related ways.

First: several phenomena that are routinely presented as characteristically
relativistic or characteristically quantum-mechanical --- an underdetermination
of physical ontology by the observable facts, an ambiguity in the passage from
classical description to quantized theory, a geometric ``phase'' with no
classical counterpart, an apparent conflict between microscopic reversibility
and macroscopic causal, dissipative behavior --- already occur inside classical
mechanics, before relativity or quantization ever enter the discussion. When
such a paradox survives the excision of everything relativistic or quantum, its
source cannot honestly be attributed to relativity or to quantum theory.

Second: even restricted to the narrow task of serving as the base for
relativistic and quantum generalization, classical mechanics is not a settled,
transparent starting point. Producing the structure that gets generalized --- a
well-defined phase space, a Hamiltonian, a consistent variational principle ---
is itself a nontrivial achievement that fails, or becomes ambiguous, for entire
classes of ordinary classical systems.

It is worth pausing on the first role (classical mechanics as a part of classical
non-relativistic theory) only long enough to insist that it is real.
The three-body problem, the precession of orbits, the stability of the solar
system, the dynamics of rigid bodies and gyroscopes, the propagation of waves in
elastic media, the onset of turbulence, and the sensitive dependence on initial
conditions studied by dynamical-systems theory are all substantial, still
partially open, areas of physics conducted entirely within the classical
framework. None of this material is a mere warm-up exercise for quantum mechanics,
and a physicist could spend a full career inside it without ever needing $\hbar$.
This book is not about that role. It is mentioned here only to make clear that
when the rest of this book treats classical mechanics as a base to be generalized,
it is not thereby claiming that this exhausts what classical mechanics is.


\section{Nonholonomic Constraints}
\label{cm:nh}

A constraint of the form $f(q,t)=0$, where $q=(q^1,\dots,q^n)$ are
generalized coordinates, is called \emph{holonomic}: it can be
solved, at least locally, to eliminate one coordinate outright, and
it reduces the dimension of the configuration space itself. A more
general class of constraints is linear in the velocities,
\begin{equation}
  \sum_{i=1}^n a_i(q,t) \dot q^i + a_t(q,t) = 0,
  \label{pfaffian}
\end{equation}
a so-called Pfaffian constraint. If the left-hand side of
\eqref{pfaffian} happens to be the total time derivative of some
function $f(q,t)$, the constraint is secretly holonomic in disguise.
If it is not --- if no such $f$ exists --- the constraint is called
\emph{nonholonomic}: it restricts the \emph{velocities} available at
each point of $Q$, without restricting which \emph{configurations}
are reachable.

The canonical example is a coin (or disk) of radius $R$ rolling
without slipping on a plane, with contact point $(x,y)$, heading
angle $\phi$, and roll angle $\theta$ as coordinates. The
rolling condition reads
\begin{equation*}
  \dot x = R\,\dot\theta\cos\phi ,\qquad
  \dot y = R\,\dot\theta\sin\phi ,
\end{equation*}
which is not integrable: it constrains the instantaneous velocity to
a two-dimensional subspace of the four-dimensional tangent space at
each point, and yet, by alternating steering and rolling, the coin
can be driven from \emph{any} configuration $(x,y,\phi)$ to any
other. The instantaneous constraint is genuine, but it places no
restriction whatsoever on the reachable set of configurations --- an
instance of what control theorists call bracket-generating, or
Chow--Rashevski\u{i}, controllability. This already runs against a
piece of intuition carried over uncritically from the holonomic
case, where a nontrivial velocity constraint always shrinks the
accessible configuration space too.

A second standard example, the \emph{Chaplygin sleigh}, makes an
even sharper point.

It is a rigid body sliding frictionlessly on a
plane, with a knife edge that permits motion only along its own
length (no sideways slip). The system conserves energy exactly ---
there is no friction anywhere in the model --- and yet its reduced
dynamics on the velocity variables is attracted to a fixed
equilibrium: left alone, the sleigh asymptotically settles into
straight-line motion with no rotation, knife edge trailing the
center of mass~\cite{chaplygin}. This
happens because the reduced equations of motion on the nonholonomic
constraint surface, unlike those of an ordinary holonomic mechanical
system, do not in general preserve phase-space (Liouville) volume,
even though they do preserve energy. A perfectly conservative
mechanical system can therefore display genuinely dissipative-looking
long-time behaviour --- attracting fixed points, basins of
attraction --- with no friction anywhere in the equations. Energy
conservation, by itself, is not enough to rule out the qualitative
signatures usually associated with dissipation once constraints stop
being holonomic.

For holonomic constraints, two routes to the equations of motion ---
D'Alembert's principle (constraint forces do no virtual work, for
virtual displacements compatible with the instantaneous constraint)
and Hamilton's principle (the true path extremizes the action among
\emph{all} paths satisfying the constraint) --- are provably
equivalent. For nonholonomic (non-integrable, velocity-linear)
constraints, this equivalence simply fails, and the failure is not a
minor technicality.

There are two natural-looking ways to generalize the variational
treatment to a system with a nonholonomic constraint of the
form~\eqref{pfaffian}.

\begin{description}
  \item[Nonholonomic (Lagrange--d'Alembert, or Chetaev) mechanics.]
  Impose the constraint directly on virtual displacements at each
  instant, exactly as in the holonomic case, and only afterwards
  derive equations of motion:
  \begin{equation}
    \frac{d}{dt}\frac{\partial L}{\partial \dot q^i}
    - \frac{\partial L}{\partial q^i}
    \;=\; \sum_a \lambda_a\, a_i^a(q),
    \qquad
    \sum_i a_i^a(q)\,\dot q^i + a_t^a(q) = 0 .
    \label{lagrange-dalembert}
  \end{equation}
  \item[Vakonomic mechanics.] (The name --- "variational axiomatic
  kind" --- was coined by Kozlov~\cite{kozlov}.)
  Instead, treat the nonholonomic constraint as a genuine constraint
  on the space of \emph{paths} from the outset, as one would an
  isoperimetric constraint in the calculus of variations, and only
  then extremize the action over that constrained set of paths.
\end{description}

For holonomic constraints these two prescriptions coincide. For
genuinely nonholonomic constraints they generically do not: it has
been shown, for instance, that for the rolling coin on an inclined
plane no nontrivial vakonomic solution reproduces the
Lagrange--d'Alembert solution at
all~\cite{zampieri}. This is not merely a
mathematical curiosity to be settled by taste: Lewis and Murray
built and instrumented a physical system --- a ball rolling without
slipping on a rotating table --- specifically to adjudicate between
the two theories, and found that the observed motion matches the
Lagrange--d'Alembert (nonholonomic) prediction; they were unable to
produce vakonomic simulations resembling the experimental data at
all~\cite{lewis+murray}. The variational principle that looks like
the most direct generalization of Hamilton's principle to velocity
constraints is, for constraints realized by an actual rolling or
sliding contact, simply the wrong physical theory.

This does not make vakonomic mechanics useless --- it correctly
describes optimal-control and sub-Riemannian-geodesic problems, where
the "constraint" genuinely is a restriction on the class of paths
being optimized over, rather than a constraint enforced
instant-by-instant by a physical reaction force. The moral is
subtler and more interesting than "one theory is right and the other
wrong": the bare constraint equation~\eqref{pfaffian}, by itself,
does not determine the dynamics. The correct equations of motion
depend on how the constraint is physically \emph{realized} ---
information that is invisible at the level of the constraint
equation alone, and that a first course in mechanics, built entirely
on holonomic examples where this ambiguity cannot arise, gives no
reason to expect.

\section{Bodies That Change Shape}
\label{cm:shape}

Drop a cat, upside down, from rest. In free fall there is no
external torque about the center of mass, so angular momentum is
conserved --- and it starts, and remains, at zero. Naive intuition
says that a body with identically zero angular momentum cannot
rotate. And yet the cat lands on its feet, having reoriented itself
by roughly $180^\circ$. The naive inference silently assumes a
\emph{rigid} body; it simply does not apply to a body that can
deform itself internally, and the cat's righting maneuver --- a
sequence of twists of its front and back halves relative to one
another --- is a real, controlled deformation, not a violation of
any conservation law.

The clean way to see why this works, due to Shapere and
Wilczek~\cite{shapere+wilczek} and to Montgomery's later "gauge
theory of the falling cat"~\cite{montgomery}, is to separate a
deformable body's configuration into its \emph{shape} $s$ (the
internal, orientation-independent degrees of freedom --- relative
joint angles, say) and an overall orientation $R\in SO(3)$ that fixes
that shape in space. Locally, the full configuration space is a
principal fiber bundle with structure group $SO(3)$ over shape space,
and the requirement that total angular momentum stay identically
zero throughout the motion is a linear relation between the shape
velocity $\dot s$ and the angular velocity $\omega$ of the body,
\begin{equation}
  \omega \;=\; -\,\mathcal A(s)\,\dot s ,
  \label{connection}
\end{equation}
where $\mathcal A(s)$ is built from the shape-dependent inertia
tensor and its couplings. This is precisely the data of a connection
(a gauge potential) on the bundle: given any path $s(t)$ in shape
space, \eqref{connection} lifts it uniquely to a path in the full
configuration space, exactly as a connection lifts a path in the
base of a fiber bundle to a horizontal path upstairs.

The consequence that matters here is that if the shape returns
exactly to where it started after a closed loop in shape space,
$s(T) = s(0)$, the orientation need \emph{not} return to where it
started. The net rotation accumulated around the loop is the
\emph{holonomy} of the connection~\eqref{connection}, and for a
small loop enclosing area $\Sigma$ it is governed, Stokes'-theorem
style, by the curvature $\mathcal F = d\mathcal A + \mathcal
A\wedge\mathcal A$ of that connection:
\begin{equation}
  \Delta R \;\approx\; \exp\!\left[\oint_{\partial\Sigma}\mathcal
  A\right] \;\approx\; \exp\!\left[\int_{\Sigma}\mathcal F\right].
  \label{holonomy}
\end{equation}
This is exactly the structure of a non-Abelian geometric (Berry,
Wilczek--Zee) phase from quantum mechanics, and of Aharonov--Bohm
holonomy in gauge theory --- reproduced in full by an entirely
classical, entirely mechanical system. It is also, viewed from the
vantage point of Section~\ref{cm:nh}, nothing but another
instance of a nonholonomic constraint: "total angular momentum
$\equiv 0$" is a Pfaffian relation between $\dot s$ and $\omega$ that
fails to be integrable to any function of shape and orientation
alone \emph{exactly when} the connection~\eqref{connection} has
nonzero curvature. The falling-cat effect and the rolling coin are,
formally, the same phenomenon in different clothing.

The same structure governs self-propulsion in a highly viscous
fluid, where inertia is negligible and net displacement over one
cycle of body deformation is possible only if the cycle of shapes is
genuinely non-reciprocal. This is Purcell's scallop theorem: a
scallop that opens and closes through the \emph{same} sequence of
shapes, forwards then backwards, achieves zero net displacement,
because time-reversing the (inertia-free) equations of motion
exactly reverses the trajectory. Net swimming is again a
holonomy/geometric-phase effect in an appropriately defined shape
bundle~\cite{shapere+wilczek-swim}.

Perhaps the most striking illustration of the underlying idea is Jack
Wisdom's result that a quasi-rigid body executing a cyclic sequence
of internal shape changes in a curved \emph{spacetime} --- as
prescribed by general relativity --- can achieve a net translation
even in complete vacuum, with nothing to push against at
all~\cite{wisdom}. The effect scales with the local spacetime
curvature and with the extent of the body's internal motion; it is
minuscule for any realistic body and is not a practical propulsion
mechanism, but as a matter of principle it is "swimming" through
spacetime using exactly the same holonomy mechanism as the falling
cat, now realized by the genuine nonlinearity of general relativity
rather than by a gauge potential on a shape bundle constructed by
hand. Once again, an outcome that sounds like it must violate
momentum conservation turns out to be perfectly consistent with
conservation laws, provided one tracks correctly what "conservation"
does and does not forbid for a body that is allowed to change shape.

\section{Open Systems}
\label{cm:open}

Ordinary Newtonian mechanics implicitly assumes a fixed set of
particles. Rockets, raindrops growing by accretion, chains lifted
off a pile, and avalanches all violate that assumption: the mass $m$
of the body under study is itself a function of time. The correct
generalization of Newton's second law to such an \emph{open system}
--- worked out around the turn of the twentieth century by
I.\,V.\ Meshchersky --- is
\begin{equation}
  m\,\frac{d\mathbf v}{dt} \;=\; \mathbf F_{\mathrm{ext}}
  + \mathbf u\,\frac{dm}{dt} ,
  \label{meshchersky}
\end{equation}
where $\mathbf u$ is the velocity of the incoming or outgoing mass
\emph{relative to the body itself}, not relative to the lab frame.
Setting $\mathbf F_{\mathrm{ext}}=0$ and integrating
\eqref{meshchersky} with $\mathbf u$ constant and antiparallel to
the direction of travel yields the Tsiolkovsky rocket equation,
\begin{equation}
  v(t) - v(0) \;=\; u\,\ln\frac{m(0)}{m(t)} .
  \label{tsiolkovsky}
\end{equation}
 The
term $\mathbf u\,\dot m$ in \eqref{meshchersky} looks, from the
point of view of the rocket body alone, exactly like an applied
force --- "thrust" --- and is routinely called one. But viewed as
part of the larger \emph{closed} system consisting of the rocket
\emph{together with} every parcel of exhaust it has ever ejected,
whose total momentum is exactly conserved, there is no force there
at all: only the ordinary transfer of momentum between two pieces of
one conserved whole. What looks like a force is, once again, an
artifact of insisting on an equation of motion for part of a system
rather than for the whole of it.

This is also a genuine pedagogical trap, worth flagging explicitly.
It is tempting to write $\mathbf F = d(m\mathbf v)/dt = m\dot{\mathbf
v} + \mathbf v\,\dot m$ and treat this directly as the equation of
motion of the open body. That expression is correct as a statement
about the rate of change of the momentum of the body \emph{alone},
but it is not, in general, equivalent to \eqref{meshchersky} and
gives the wrong dynamics unless the accreted or ejected mass happens
to leave (or join) with zero velocity relative to the body, i.e.\
$\mathbf u \equiv \mathbf v$ --- a special case, not the generic one,
and not the one realized by a rocket engine, whose whole purpose is
to eject mass at high relative velocity.


A second, quite different kind of "openness" arises whenever a
system exchanges energy with an environment that is not explicitly
modelled --- friction, drag, radiation reaction. None of these
forces derive from a potential in the ordinary sense, so the
Lagrangian and Hamiltonian formalisms, built for closed and
energy-conserving systems, do not apply directly. Two classical
devices repair this.

The first is Rayleigh's dissipation function $\mathcal
F_{\mathrm{diss}}(\dot q)$, quadratic in the generalized velocities
for linear damping, which modifies the Euler--Lagrange equations to
\begin{equation}
  \frac{d}{dt}\frac{\partial L}{\partial \dot q^i}
  - \frac{\partial L}{\partial q^i}
  \;=\; -\,\frac{\partial \mathcal F_{\mathrm{diss}}}{\partial \dot
  q^i}.
  \label{rayleigh}
\end{equation}
This correctly reproduces the (irreversible) equations of motion,
but it is a hybrid: $\mathcal F_{\mathrm{diss}}$ is not itself part
of a genuine, self-contained stationary-action principle for $L$
alone.

The second, due to Bateman (1931), is more radical and more
illuminating: a fully conservative, symplectic Lagrangian \emph{can}
be written for a damped oscillator, but only at the price of
doubling the number of degrees of freedom. Introduce an auxiliary
"mirror" coordinate $y$ alongside the physical coordinate $x$,
\begin{equation}
  L_{\mathrm{dual}}(x,\dot x,y,\dot y) \;=\;
  m\,\dot x\,\dot y \;+\; \frac{\gamma}{2}\bigl(x\dot y - \dot x
  y\bigr) \;-\; k\,xy ,
  \label{bateman}
\end{equation}
whose Euler--Lagrange equations are precisely the damped oscillator,
$\ddot x + (\gamma/m)\dot x + \omega_0^2 x = 0$, for the physical
coordinate $x$, paired with a time-reversed, exponentially
\emph{amplified} equation for $y$. The energy the
physical oscillator loses is exactly the energy its mirror partner
gains, so the doubled system is conservative even though either half
alone is not. This is a small, entirely classical and elementary
demonstration of an idea that later became central to the theory of
open quantum systems: dissipation and irreversibility in a subsystem
are not exceptions to an underlying conservative dynamics, but the
generic appearance of a conservative dynamics once part of the full
system --- here, quite literally, the mirror half --- has been
discarded from the description.


Not every attempt to write a closed-looking equation for an open
system ends benignly. An accelerating point charge radiates, and by
momentum conservation it must experience a recoil, or self-force, in
consequence. Treating the charge \emph{in isolation} --- an open
system in the strict sense, since the charge together with its own
radiated field is really a system with infinitely many degrees of
freedom, formally "integrated out" down to a single effective force
term --- leads to the (nonrelativistic) Abraham--Lorentz equation,
\begin{equation}
  m\,\ddot{\mathbf x} \;=\; \mathbf F_{\mathrm{ext}} \;+\;
  m\tau\,\dddot{\mathbf x} ,
  \qquad
  \tau \;=\; \frac{q^2}{6\pi\varepsilon_0\,m\,c^3}
  \;=\; \frac{2}{3}\,\frac{r_e}{c},
  \label{abraham-lorentz}
\end{equation}
where $r_e$ is the classical electron radius. Because
\eqref{abraham-lorentz} is third order in position, it requires
more initial data than position and velocity to determine a unique
solution, and generic initial data produce \emph{runaway} solutions
in which the self-acceleration grows without bound even with
$\mathbf F_{\mathrm{ext}}=0$. Suppressing the runaway by hand forces
the opposite pathology: the charge must begin accelerating
\emph{before} an external force is switched on --- a small, short-lived,
but genuinely acausal-looking pre-acceleration. (In practice, both
pathologies signal that \eqref{abraham-lorentz} is being pushed
outside its domain of validity, on timescales comparable to $\tau
\sim 10^{-24}\,\mathrm s$; treated properly, as an effective equation
valid only for slowly varying external forces --- via, e.g., the
Landau--Lifshitz reduction of order --- the pathologies disappear.)
The lesson is a cautionary complement to the rest of this section:
"integrating out" the unmodelled part of an open system is not
always harmless bookkeeping. Sometimes it is, as with the Meshchersky
equation; sometimes, as here, it produces an equation that is simply
malformed unless handled with real care.

\section{Simple Systems}
\label{cm:ss}

Let us restrict ourselves with Conservative Holonomic Systems or
more specifically with systems consisting of finite number of
interacting point particles.
There is a great number of books on the subject
(see e.g\cite{Arnold-cm},\cite{Landau-Lifshitz-1},
\cite{Lanczos},\cite{Goldstein}). 

In Ideal world of classical mechanics space is three-dimensional  and euclidean
$\mathbb{R}^3$, mass is real number and time is one-dimensional $\mathbb{T}$.
So universe of classical mechanics is $\mathbb{T} \times \mathbb{R}^3$. 

 The  group of galilean transformations acts on universe \cite{Arnold-cm}. 
Every galilean transformation can be written in a unique way as the composition of a rotation, a translation,
and a uniform motion ($g = g_1 \circ g_2 \circ g_3$), where
\[g_1(t,\bm x)=(t,\bm x+\bm v t)\qquad \forall t\in \mathbb{R},\bm x \in\mathbb{R}^3\] 
is uniform motion with velocity $\bm v \in\mathbb{R}^3$,
\[g_2(t,\bm x)=(t+t_0,\bm x+\bm x_0)\qquad \forall t\in \mathbb{R},\bm x \in\mathbb{R}^3\] 
is translation of the origin, and
\[g_3(t,\bm x)=(t,O\bm x)\qquad \forall t\in \mathbb{R},\bm x \in\mathbb{R}^3\] 
rotation of the coordinate axes ($O:\mathbb{R}^3 \rightarrow \mathbb{R}^3$ is an orthogonal transformation).

Let $M$ be a set. A one-to-one correspondence 
\[\phi_1: M \rightarrow \mathbb{R} \times \mathbb{R}^3\]
is called a galilean (inertial) coordinate system on the set $M$. A coordinate system $\phi_2$ moves
uniformly with respect to $\phi_1$  if $\phi_1\circ\phi_2^{-1}$ is a galilean
transformation. 

Inertial coordinate systems  possess the following two properties:
\begin{enumerate}
    \item All the laws of nature at all moments of time are the same in all inertial coordinate systems;
    \item  All coordinate systems in uniform rectilinear motion with respect to an inertial one are themselves inertial.
\end{enumerate}
This is called `Galileo's principle of relativity'.

The simplest object of classical mechanics is point mass (particle) i.e. object whose shape, 
sizes and internal structure are inessential, so mass (positive scalar $m$) is the only parameter. 
The state of a particle is fully determined by position ($\bm x \in\mathbb{R}^3$) 
and velocity ($\bm v \in\mathbb{R}^3$). 
More accurately: A motion in $\mathbb{R}^n$ is a differentiable mapping 
\[\bm x: \mathbb{I}\rightarrow \mathbb{R}^3, \mathbb{I} \subset \mathbb{R}.\] 
The derivative 
\[\dot {\bm x}(t)=\frac{d{\bm x}}{dt}(t)\]
is called the velocity vector. 
The derivative 
\[\ddot {\bm x}(t)=\frac{ d^2 {\bm x}}{dt^2}(t)\]
is called the acceleration vector. 
The image of motion (curve in $\mathbb{R}^3$) called trajectory and  curve in 
$\mathbb{T} \times \mathbb{R}^3$ which appears  as the graph of a motion, is called a world line.

The initial state of a mechanical system (the totality of positions and
velocities of its points at some moment of time) uniquely determines all of
its motion. In particular, the initial positions and velocities determine the accelerations.
This is called `Newton's principle of determinacy'.

\subsection{Newtonian mechanics}

Let us consider the mechanical system consisting of $n$ particles 
with masses $m_j,j=1\ldots n$ and word lines $\bm x_j(t)$.
Suppose that $n$ vector functions   
\[{\bm F}_j:\mathbb{R}^{3 n}\times\mathbb{R}^{3 n}
\times\mathbb{R}\rightarrow\mathbb{R}^3,j=1\ldots n\]
are defined. Then the motion of all particles is determined by equations
\begin{equation}
  \label {newtoneq}
    m_j \ddot {\bm x}_j={\bm F}_j({\bm x}_k,\dot {\bm x}_k,t) \quad j,k=1\ldots n 
\end{equation}
and initial conditions
\begin{equation}
\label {newtonic}
    {\bm x}_j(t_0)={\bm y}_j\quad\dot {\bm x}_j(t_0)={\bm v}_j  \quad j=1\ldots n 
\end{equation}

Newton's equation \eqref{newtoneq} is  the basis of mechanics. Functions 
$\bm F_j$ are called forces they should be determined  for each specific 
mechanical system. If our system is isolated i.e. includes explicitly all 
interacting particles, then galilean invariance puts some restrictions 
on forces.

Namely 
\begin{equation}
  \label {newtoneq-isolated}
    m_j \ddot {\bm x}_j=
    {\bm F}_j({\bm x}_j-{\bm x}_k,\dot {\bm x}_j-\dot {\bm x}_k)\quad j,k=1\ldots n 
\end{equation}
practically such systems occur in celestial mechanics. In most other cases the
system considered is not isolated and forces of general kind \eqref{newtoneq} 
appear as a result of interaction with environment.

Mathematics of  Cauchy problem (\ref{newtoneq},\ref{newtonic}) 
one can find in e.g. \cite{Arnold-ode}. Generally the unique solution exists if
all functions are smooth enough. However smoothness may be violated in 
physically important cases, anyway functions $\bm x_j(t)$ should be continuous,
$\bm v_j(t)$  may have finite jumps as a results of strong force acting for 
very short time (e.g. elastic collision).

Note that initial conditions have nothing to do with causality. Causality makes
sense in Real world only, in Ideal world we are free to set final conditions or
fix $\bm x_j(t_1)$ and $\bm v_j(t_1)$ for some intermediate time $t_1$. In
many cases we are interested in solution of equation \eqref{newtoneq} with
fixed initial velocities and final coordinates  (aiming) or fixed coordinates
on both ends (variational problem). Anyway when conditions are set at two
different times  solution may be not unique (unless force is changing
relatively slow) \footnote {I discovered this fact around 1980 and was very
proud of myself until I found papers \cite{Maslov-1} and \cite{Maslov-2}.}.

\paragraph{Constraints} 

In addition to Newton's equations the motion of the system may be restricted by
additional equations or inequalities. Constraints introduce two types of
difficulties in the solution of mechanical problems. First, the coordinates
$\bm x_j$ are no longer all independent, since they are connected by the
equations of constraint; hence the equations of motion  are not all
independent. Second, the forces of constraint  are not furnished a priori. They
are among the unknowns of the problem and must be obtained from the solution we
seek. Indeed, imposing constraints on the system  is simply another method of
stating that there are forces present in the problem that cannot be specified
directly but are known rather in terms of their effect on the motion of the
system. There is no general method of solving (even numerically) these
problems, they are tackled on case by case basis.

For the moment we should restrict ourselves with so called holonomic
constraints which have the form:
\begin{equation}
        f(\bm x_j,t)=0 \label {hol-const}
\end{equation}

\paragraph{Numerical solutions}

Generally Newton's equations are to be solved numerically which is a difficult
task. Two large class of methods are in use: 
\begin{enumerate} 
    \item "Usual" one-time integrators when differential equations are replaced
        with finite difference ones defining snapshots at discrete time steps.
        These integrators allow accurate description of standard problem:
        initial conditions, finite time interval. 

\item Geometric integrators which are less accurate for any single trajectory
    but accurately present qualitative picture for large set of initial data
        and long times. Such integrators may calculate  coordinates and momenta
        at different times
\end{enumerate}
One can look for details in \cite{Hairer-Lubich-Wanner} for example. 

If we allow only potential forces i.e. 
\begin{equation}
  \label {potforce}
        \bm F_j(\bm x_k,t)=-\frac{\partial U(\bm x_k,t)}{\partial \bm x_j} \quad j,k=1\ldots n 
\end{equation}
we shall see that equations of motion (\ref{newtoneq}) can be written as 
\begin{equation}
   \label {lagr}
   \begin{split}
       {\bm p}_j&=\frac{\partial L({\bm x}_k,{\bm v}_k,t)}{\partial {\bm v}_j} \\ 
        \dot {\bm p}_j&= \frac{\partial L({\bm x}_k,{\bm v}_k,t)}{\partial {\bm x}_j}\quad j,k=1\ldots n 
   \end{split}
\end{equation}  
where 
\begin{equation}
  \label{lagr3}
        L({\bm x}_k,{\bm v}_k,t)=\Sigma_1^n \frac{m_k}{2} {\bm v}_k \cdot {\bm v}_k-U({\bm x}_k,t) 
\end{equation}

If we can solve the equation 
\[ \bm p_j=\frac{\partial L(\bm x_k,\bm v_k,t)}
{\partial \bm v_j} \] 
for $ \bm v_j $
which of course is not always possible, we can write the equations \eqref{lagr}
in more symmetrical way:
\begin{equation}
   \label {ham}
   \begin{split}
       \dot {\bm x}_j&=\frac{\partial H(\bm x_k,\bm p_k,t)}{\partial \bm p_j}  \\ 
        \dot {\bm p}_j&= -\frac{\partial H(\bm x_k,\bm p_k,t)}{\partial \bm x_j}\quad j,k=1\ldots n    
   \end{split}
\end{equation}
where 
\begin{equation}
        H(\bm x_k,\bm p_k,t)=\Sigma_1^n \frac{1}{2 m_k} \bm p_k \cdot \bm p_k+U(\bm x_k,t) \label{ham3}
\end{equation}
Note that functions $L$ (Lagrangian) and $H$ (Hamiltonian) are connected by Legendre transform:
\begin{equation}
        L(\bm x_k,\bm v_k,t)=\Sigma_1^n  \bm p_k \cdot \bm v_k -H(\bm x_k,\bm p_k,t)|_{\bm p_k=m_k\bm v_k}
\quad k=1\ldots n \label {lagr-ham}
\end{equation}


\subsection{Principle of least action}

\begin{definition}
        The equation
        \[\frac{d}{d t}(\frac{\partial L}{\partial \dot {\bm x}})-\frac{\partial L}{\partial  {\bm x}}=0 \]
        is called the Euler-Lagrange equation for the functional
        \[\Phi=\int_{t_1}^{t_2} L(\bm x,\dot {\bm x},t) dt \]
\end{definition}
Let \( \gamma= \bm x(t) \quad t_1 \le t \le t_2 \) be a curve  in  $\mathbb{R} \times \mathbb{R}^n$ and 
$L: \mathbb{R}^n \times \mathbb{R}^n \times \mathbb{R}\rightarrow \mathbb{R}$ a function of $2 n +1$
variables, then
\begin{theorem}
    The curve $\gamma$ is an extremal of the functional \[\Phi=\int_{t_1}^{t_2} L(\bm x,\dot {\bm x},t) dt \]
        on the space of curves joining $(t_1, \bm x_1)$ and $(t_2,\bm x_2)$, if and only if the Euler-
Lagrange equation is satisfied along $\gamma$.
\end{theorem}
This is a system of n second-order equations, and the solution depends on
$2 n$ arbitrary constants. The $2 n$ conditions $\bm x(t_1) = \bm x_0$ and $\bm x(t_2) = \bm x_2$ are used
for finding them.
Two important remarks:
\begin{enumerate}
        \item There may be many extremals connecting two
given points or none at all.
\item The condition for a curve  to be an extremal of a functional does not depend
on the choice of coordinate system. This paves way to geometrisation of mechanics.
\end{enumerate}

It is easy to see that equation of motion (\ref{lagr}) are actually  Euler-Lagrange equations 
for lagrangian \eqref{lagr3}. 
So we have  Hamilton's principle of least action.


I believe that least action principle  is nothing but (powerful) heuristic. Indeed, to compute 
action functional we should define explicitly {\em all} curves connecting given pair of points.
This seems to be impossible or at least very difficult problem. Fortunately we do not need it.
All what we really need is to set configuration space, choose Lagrange function 
\eqref{lagr3} and get equations of motion \eqref{lagr}. It does not matter at all
if we have minimum, maximum, saddle point or whatever property of some functional.

On the other hand having equations of motion as system of ordinary differential equations,
we can transform it into equivalent one, say by multiplying every equation by arbitrary 
function which never becomes zero. This new system may correspond to another Lagrangian and
the difference between new and old one need not come down to total derivative. Most probably
new system will not be Lagrangian at all.

Let us look at two-dimensional harmonic oscillator
\footnote{This is a simple example of difference between the
symmetry of equations and that of Lagrangian, which will be discussed later}.

Usual Lagrangian is
\begin{equation}
   \label {ho1}
    L_1= \frac{1}{2} (\dot x^2+\dot y^2) -\frac{ \omega^2}{2} ( x^2+ y^2)
\end{equation}
which yields equation  of motions:
\begin{equation}
   \label {hoe1}
   \begin{split}
       \ddot x+ \omega^2 x & =0 \\
       \ddot y+ \omega^2 y & =0 
   \end{split}
\end{equation}
Same equations of motion come form 
another Lagrangian
\begin{equation}
   \label {ho2}
    L_2= \frac{1}{2} (\dot x^2-\dot y^2) -\frac{ \omega^2}{2} ( x^2- y^2)
\end{equation}
or
\begin{equation}
   \label {ho3}
    L_3= \dot x\dot y - \omega^2  x y
\end{equation}
corresponding  Hamiltonians 
\begin{equation}
   \label {ho2h}
    H_2= \frac{1}{2} ( p_x^2-p_y^2) +\frac{ \omega^2}{2} ( x^2- y^2)
\end{equation}
and
\begin{equation}
   \label {ho3h}
    H_3= p_x p_y +\omega^2  x y
\end{equation}
are not even positive definite!

\subsection{Gauge invariance} 

From now on we shall use generalized coordinates $ q=\{q_j, j=1\ldots n\} $, where $ n $ is the number of
degrees of freedom. These coordinates may have or have not physical meaning.

We are used to think that gauge invariance is about electromagnetism, actually it appears
already in classical mechanics. It is easy to prove that adding the total derivative
\begin{equation}
   \label {tot-der}
    \frac{d \Phi (q,t)}{d t}=\frac{\partial \Phi(q,t)}{\partial t} + \frac{ \partial \Phi (q,t)}{\partial q_j} \dot q_j
\end{equation}
to Lagrangian does not change equations of motion. However canonical momenta
\begin{equation}
   \label {can-mom}
   p_j=\frac{ \partial L(q,\dot q,t)}{\partial \dot q_j}
\end{equation}
becomes \[ P_j=p_j+\frac{\partial \Phi}{\partial q_j} \] Transformation $ p,q \rightarrow P,q $ is evidently
canonical and new Hamiltonian becomes 
\[ H_{new}=H(p(P),q)- \frac{\partial \Phi}{\partial t}\]
The fact that canonical momenta have no physical meaning is not a weak but strong point of Hamiltonian 
mechanics.

I skip the  discussion of standard things like Poisson brackets, integrals of motion and all that.

\subsection{"Spin" in constant magnetic field}

We used to natural systems where hamiltonian is quadratic form over momenta,
but there are physical systems with hamiltonians linear over momenta.
So lagrangian does not exist. Let us consider 
dynamic of point magnetic moment in constant magnetic field has not only
clear physical meaning but also significant practical value. I investigated
this problem (actually slightly more general one) in connection with control of MRI
scanner \cite{RPZ4},\cite{RPZ3},\cite{RPZ2},\cite{RPZ1}.


Hamiltonian is $ H=\bm M \cdot \bm B $ where $\bm M$ is magnetic moment 
and $ \bm B $ is constant magnetic field.
It is convenient to write $ \bm M=\frac{e}{m c}\bm S $ where $ \bm S $
is angular moment ("spin")  and  include gyromagnetic ratio into magnetic field
$ \bm \Omega = \frac{e}{m c}\bm B$. Components of $ \bm S $ are not
canonical variables, their Poisson brackets are
\begin{equation}
   \label {pb}
    \{S_i,S_j\}= \epsilon_{ijk} S_k
\end{equation}
However they define phase space (sphere)  and equation of motion:
\begin{equation}
   \label {sp1}
    \dot S_i=\{S_i,H\}
\end{equation} 
or 
\begin{equation}
   \label {sp2}
    \dot {\bm S} =-\bm S \times \bm \Omega
\end{equation}

In canonical variables 
\begin{equation}
   \label {sp3}
   \begin{split}
       S_x & = \sqrt{S^2-p^2} \cos q \\
       S_y & = \sqrt{S^2-p^2} \sin q \\
       S_z & = p,\quad -S\le p \le S
   \end{split}
\end{equation}
Let $ \bm \Omega =(0,0, \Omega) $ then $ H=p \Omega $ and canonical equations are:
\begin{equation}
   \label {sp4}
   \begin{split}
       \dot p & =0\\
       \dot q & = \Omega
   \end{split}
\end{equation}
Lagrangian $ L=0 $, so we can not get equations of motion as second order equation.

Moment $ p $ is integral of motion which is fine but velocity is not connected to 
moment at all, it is constant. For pair of points in phase space generally there is
no trajectory passable for given time. Variation problem simply can not be posed.

Hamiltonian $ H=p \Omega $ looks like that for harmonic oscillator in
action-angle variables and so we can  transform it to usual quadratic form, but
the condition $  -S\le p \le S$ makes this transformation pure formal.

\section{EOM: Conditions and Problems}
\label{cm:eom}

This section organizes the standard problem types associated with an
equation of motion (EOM) along two axes: (A) \emph{what data specifies
the problem} (the conditions imposed), and (B) \emph{what feature of
the solution's structure is being investigated}, once the EOM itself
is fixed.

\subsection{A. By the Conditions Imposed}

\paragraph{IVP / Cauchy problem}
  \textbf{Condition:} the full state (position and velocity, or a
  phase-space point) at one instant $t_0$.
   \textbf{Governing results:} the Picard--Lindel\"of theorem gives
  local existence and uniqueness for Lipschitz forces; global existence
  requires a no-blow-up bound.
   This is what time-stepping integrators (RK4, symplectic
  schemes, etc.) actually solve.

\paragraph{Boundary Value Problem (BVP)}
   \textbf{Condition:} data at two or more distinct times/points,
  e.g. $x(t_0) = x_0$, $x(t_1) = x_1$ --- not a full state at one
  instant.
   Existence and uniqueness are \emph{not} automatic (unlike the
  IVP): zero, one, or many solutions can occur.
   \textbf{Prototypes:} the brachistochrone problem, geodesics
  between two points, Lambert's problem, and any Euler--Lagrange
  equation arising from a fixed-endpoint action principle.
   \textbf{Numerically:} shooting methods, relaxation /
  finite-difference methods, collocation.

\paragraph{Eigenvalue Problem (a BVP subtype)}
   A linear EOM together with boundary conditions admits
  nontrivial solutions only for special parameter values ---
  normal-mode frequencies, quantized energies.

\paragraph{Scattering Problem}
   \textbf{Condition:} an asymptotic free state as $t \to
  -\infty$ (or $r \to \infty$).
   \textbf{Sought:} the asymptotic outgoing state as $t \to
  +\infty$ --- the $S$-matrix or cross-section --- given a localized
  interaction region in between.
   Differs in kind from the IVP/BVP: the conditions sit at
  infinity, and often the trajectory itself is discarded --- only the
  in $\to$ out map matters.

\paragraph{Periodicity Condition}
   Search for $x(t+T) = x(t)$. Effectively a BVP with the two
  endpoints identified, but $T$ is itself unknown, making it a
  nonlinear eigenvalue-type problem.
   \textbf{Poincar\'e--Bendixson theorem} (2D flows): a bounded
  trajectory that does not approach a fixed point must approach a
  periodic orbit.

\paragraph{Inverse Problem}
   \textbf{Condition:} observed trajectories or scattering data.
   \textbf{Sought:} the force law / potential / EOM itself
  (inverse scattering; Newton inferring $1/r^2$ from Kepler's laws is
  the archetype).

\paragraph{Control / Optimization Problem}
   \textbf{Sought:} an external input, or a member of a family of
  EOMs, that steers the system to a target state while extremizing a
  cost functional --- Pontryagin's maximum principle.
   Structurally a BVP with an added adjoint (costate) system.

\subsection{B. By the Question Asked About Solution Structure}
Given a fixed EOM, the following ask about the qualitative or global
structure of its solutions rather than about how the problem was
posed.

\begin{center}
\begin{tabular}{@{}p{4cm}p{5cm}p{6cm}@{}}
\hline
\textbf{Problem} & \textbf{What it asks} & \textbf{Key tools / results} \\
\hline
Existence / uniqueness / well-posedness &
Does a solution exist, is it unique, does it depend continuously on the data? &
Picard--Lindel\"of; Hadamard well-posedness \\
Stationary solutions &
Fixed points $\dot{x} = f(x) = 0$, or static critical points $\nabla V = 0$ &
--- \\
Linear stability &
Behavior under infinitesimal perturbation of a fixed point / orbit &
Jacobian eigenvalues (node / saddle / focus / center); Floquet multipliers for periodic orbits \\
Nonlinear (Lyapunov) stability &
Does the full nonlinear flow stay near, or converge to, equilibrium? &
Lyapunov functions; LaSalle invariance principle \\
Structural stability &
Does the qualitative phase portrait survive small changes to the EOM itself? &
Leads into bifurcation theory \\
Bifurcation &
Qualitative change in fixed points / orbits as a parameter varies &
Saddle-node, pitchfork, transcritical, Hopf, period-doubling \\
Attractors \& basins &
Long-time ($t \to \infty$) fate of trajectories in dissipative systems &
Fixed point / limit cycle / torus / strange attractor; basin of attraction \\
Chaos &
Sensitive dependence on initial conditions &
Lyapunov exponents; Poincar\'e sections; homoclinic tangles (Melnikov method); routes to chaos (period-doubling, Ruelle--Takens, intermittency) \\
Integrability &
Enough conserved quantities in involution to reduce the problem to quadratures &
Liouville integrability; action-angle variables; Noether symmetries; Lax pairs; Painlev\'e test \\
Perturbation theory &
Fate of near-integrable systems under small perturbation &
Regular / singular perturbation; multiple scales; averaging; KAM theorem \\
\hline
\end{tabular}
\end{center}

\subsection{Structural Notes}
   \textbf{Conservative (Hamiltonian) systems have no attractors}
  --- Liouville's theorem forbids phase-space contraction --- so their
  long-time behavior is governed instead by the
  integrability/KAM/chaos axis (invariant tori surviving or breaking
  up into chaotic layers), not by attractor classification.

   \textbf{Dissipative systems} are where ``attractor'' language
  belongs; stability and bifurcation theory there describe how the
  attractor set itself changes.

   \textbf{Scattering is best read as an asymptotic BVP}: the
  conditions sit at $t = \pm\infty$ rather than at finite times, which
  is exactly why it needs its own machinery ($S$-matrix, asymptotic
  completeness) rather than shooting/relaxation methods.

   Bifurcation theory is the natural bridge between ``stability
  of a single solution'' and ``existence of qualitatively new families
  of solutions'' (e.g., a fixed point losing stability via a Hopf
  bifurcation and handing off to a limit cycle, which can itself
  period-double into chaos).

